\begin{frame}[fragile]{What is Functional Programming (FP)}

\begin{itemize}
	\item In the OOP paradigm, we have two different entities in programming: data (fields) and code (methods).
	\item We treat data and code differently in OOP, as in this example Dog class.  
\end{itemize}

\begin{Code}{}%
\begin{lstlisting}[firstnumber=1, label=glabels, xleftmargin=0pt, basicstyle=\scriptsize]% 

class Dog {
  String name = "Sparky"; // data (fields)
  
  void bark() { // code (methods)
     ... 
  }
}
\end{lstlisting}%   
\end{Code}%
\end{frame}

\begin{frame}[fragile]{}

\begin{itemize}
	\item In the FP paradigm, there is no difference between the data and code; in other words, we treat the code as the same as the data.
	\item We can define a variable (line 1) and a function (line 2), but the function is nothing more than a lambda expression (line 3) as a value (data). 
\end{itemize}

\begin{Code}{}%
\begin{lstlisting}[firstnumber=1, label=glabels, xleftmargin=0pt, basicstyle=\small]% 

val x = 3
add x y = x + y
val add = (x, y) => (x + y) 
\end{lstlisting}%   
\end{Code}%
\end{frame}

\begin{frame}[fragile]{}

\begin{itemize}
	\item Using this idea, a function can be an argument to other high-level functions. 
	\item In this example, we can give add or sub function as an argument to the \verb|high_lvel_function|. 
	\item This is possible, because code (function) is the same as data (variables) in FP. 
\end{itemize}

\begin{Code}{}%
\begin{lstlisting}[firstnumber=1, label=glabels, xleftmargin=0pt, basicstyle=\scriptsize]% 

val add = (x, y) => (x + y)
val sub = (x, y) => (x - y)
val high_level_function = (f, x, y) => f(x, y)
high_level_function(add, 10, 20) // returns 10 + 20
high_level_function(sub, 10, 20) // returns 10 - 20 
\end{lstlisting}%   
\end{Code}%

\end{frame}

\begin{frame}[fragile]{JavaScript and Lisp/Scheme}

\begin{itemize}
	\item JavaScript uses functional programming approaches as it is based on Lisp/Scheme.
	\item The first method (line 1) is used to define a method in JavaScript.
	\item However, it is OK to use the FP way to treat a function as a data, like Lisp/Scheme (line 5); we call it lambda ($\lambda$) expression (lines 2 -- 3). 
\end{itemize}

\begin{Code}{}%
\begin{lstlisting}[firstnumber=1, label=glabels, xleftmargin=0pt, basicstyle=\scriptsize]% 

function add(x,y) {return x + y; }
var add2 = function (x, y){ return x + y; }
var add3 = (x, y) => {return x + y;}

(define (add) (lambda (x y) (+ x y))) ; Scheme code 
\end{lstlisting}%   
\end{Code}%
\end{frame}

\begin{frame}[fragile]{}

\begin{itemize}
	\item In Lisp/Scheme, functions are the same as data with a type (lines 3 -- 4 for function and line 2 for data). 
	\item So, any function can be used as an argument to any function (line 5). 
	\item In JavaScript, we can now understand why a function can be used as an argument (line 8) to another function (lines 11 -- 12).
\end{itemize}

\begin{Code}{}%
\begin{lstlisting}[firstnumber=1, label=glabels, xleftmargin=0pt, basicstyle=\scriptsize]% 

// scheme
(define x 3)
(define (high_order_function f) ( ;; f is a function
  (apply f '(2 3)))) ; or (f 2 3)
(high_order_function (lambda (x y) (+ x y)))

// JavaScript (f is a function)
function high_order_function(f) {return f(2, 3);}
var r = high_order_function(function (x,y) {return x+y;})
var r = high_order_function((x,y) => {return x+y;})
\end{lstlisting}%   
\end{Code}%

\end{frame}

\begin{frame}[fragile]{}
\begin{itemize}
    \item We can understand JavaScript programming better when we are familiar with functional programming styles. 
	\item This is an example of event-driven programming in JavaScript using functional programming.
	\item It is interesting to notice that Java adopts the functional programming style for its event-driven programming (CSC 360).  
\end{itemize}

\begin{Code}{}%
\begin{lstlisting}[firstnumber=1, label=glabels, xleftmargin=0pt, basicstyle=\tiny]% 

<p id="hello">Hello</p>
<button id="button">Button</button>

<script>
document.getElementById('button').addEventListener(
  'click', 
  function() { // <- function is given as an argument
   document.getElementById('hello').innerHTML = "Good bye";
  });
</script>
\end{lstlisting}%   
\end{Code}%
\end{frame}

\begin{frame}[fragile]{}

\begin{itemize}
	\item There is another example; we also can understand why there are two ways to iterate over an array. 
	\item Non-functional (procedural) way with variables and for loop (lines 1 - 3).
	\item Functional way to use function (lambda expression) as an argument of a function (lines 5 - 7).
	\item Java also adopts this functional programming style (CSC 360). 
\end{itemize}

\begin{Code}{}%
\begin{lstlisting}[firstnumber=1, label=glabels, xleftmargin=0pt, basicstyle=\scriptsize]% 

var x = [1,2,3,4];
for (var i = 0; i < x.length; i++) {
  console.log(x[i])
}
x.forEach(function(a) { // this is better
  console.log(a)
});
\end{lstlisting}%   
\end{Code}%
\end{frame}
